% =====================================================================
% Boutique Pipeline — User Guide (image-first rewrite)
%
% Audience: an operator who launches the Streamlit wizard exposed by
% src/pipelines/boutique_pipeline.py and wants to turn a phone recording
% into a per-ride PDF report. The body of this section is the
% \emph{quick-start basics}: a 2-minute walkthrough plus a one-page
% visual cheat-sheet for "is this segment good or bad". Everything that
% does not affect the operator's day-to-day workflow has been moved into
% the \emph{Advanced} subsection at the end of this section, which the
% casual user can ignore.
% =====================================================================

\section{Boutique Pipeline: A User's Guide}
\label{sec:boutique-pipeline-guide}

This section is the ready-set-go manual for the Streamlit app at
\texttt{src/pipelines/boutique\_pipeline.py}. It is intentionally
image-first: each step shows you what the screen actually looks like,
and each "what to look for" check has a side-by-side good vs.\ bad
example so you can decide in seconds. If you only want to ship a
report, read \S\ref{sec:boutique-quickstart} and then
\S\ref{sec:boutique-good-vs-bad}; everything after that is optional.

\subsection{Quick-start: five clicks to a report}
\label{sec:boutique-quickstart}

\paragraph{Launch.} From the project root, with the venv activated:

\begin{verbatim}
venv/bin/python -m streamlit run src/pipelines/boutique_pipeline.py
\end{verbatim}

\noindent Streamlit prints a localhost URL (default
\texttt{http://localhost:8501}). The wizard advances with the primary
button at the bottom-right of each page. Below is the entire happy-path
workflow.

% ---------- Step 1 ----------
\begin{figure}[H]
\centering
\includegraphics[width=0.95\linewidth]{boutique/ui_step1_landing.png}
\caption{\textbf{Step~1 --- Landing.} One paragraph of orientation;
click \emph{Start \textrightarrow} to begin.}
\label{fig:bq-step1}
\end{figure}

% ---------- Step 2 ----------
\begin{figure}[H]
\centering
\begin{minipage}[t]{0.49\linewidth}
\centering
\includegraphics[width=\linewidth]{boutique/ui_step2_picker.png}
\par\textbf{2a.} Pick \emph{Use phone DB} (in-house) or
\emph{Upload a file} (everyone else).
\end{minipage}
\hfill
\begin{minipage}[t]{0.49\linewidth}
\centering
\includegraphics[width=\linewidth]{boutique/ui_step2_upload_form.png}
\par\textbf{2b.} For uploads: a CSV/XLSX with a \emph{time} column
(seconds or ms; auto-detected) and a \emph{vertical-acceleration}
column (m/s$^2$). The first row must be a header.
\end{minipage}
\caption{\textbf{Step~2 --- Data.} Two input modes. Switch back and
forth at any time using \emph{change input source} in the upper-right
corner.}
\label{fig:bq-step2}
\end{figure}

% ---------- Step 3 ----------
\begin{figure}[H]
\centering
\includegraphics[width=0.97\linewidth]{boutique/panels/clean_step3_overview__signal.png}
\caption{\textbf{Step~3 --- Segmentation.} The detector runs
automatically. Each detected ride is shaded
(\textcolor{blue!70!black}{blue $=$ up},
\textcolor{purple!60!red}{magenta $=$ down}); the currently selected
segment is \textcolor{orange!90!black}{orange}. If the segments look
right, click \emph{Predict \textrightarrow}. If something is off, see
\S\ref{sec:boutique-good-vs-bad} for what to look for and
\S\ref{sec:boutique-fixing} for how to fix it.}
\label{fig:bq-step3}
\end{figure}

% ---------- Step 4 ----------
\begin{figure}[H]
\centering
\includegraphics[width=0.95\linewidth]{boutique/crops/clean_step4_metrics.png}
\caption{\textbf{Step~4 --- Prediction.} Both algorithms (Trapezoid
pulse-pair, in blue; ZUPT, in green) run automatically and one bar per
algorithm per segment is plotted. Positive bars mean the cabin went
\emph{up}. The two bars should be similar height; large disagreement
flags a noisy segment (\S\ref{sec:boutique-advanced}). When you are
happy, click \emph{Generate report \textrightarrow}.}
\label{fig:bq-step4}
\end{figure}

% ---------- Step 5 ----------
\begin{figure}[H]
\centering
\includegraphics[width=0.85\linewidth]{boutique/crops/clean_step5_top.png}
\caption{\textbf{Step~5 --- Report.} An inline PDF preview is rendered
in the page. Click \emph{Download Summary PDF} to save the Hebrew
multi-page report (one page per ride; cover page with the per-ride
table). That's the entire pipeline.}
\label{fig:bq-step5}
\end{figure}

\paragraph{One sentence per step.}
\begin{enumerate}
  \item Read the orientation, click \emph{Start}.
  \item Pick a source, click \emph{Next}.
  \item \emph{Glance} at the segments (next subsection); if good,
    click \emph{Predict}.
  \item Glance at the bar chart; if both algorithms agree, click
    \emph{Generate report}.
  \item Click \emph{Download}. Done.
\end{enumerate}

The \emph{only} step that needs your attention is step~3. The next
subsection is the visual cheat-sheet for it.

\subsection{Step 3 cheat-sheet: is this segment good or bad?}
\label{sec:boutique-good-vs-bad}

The segmentation page has four diagnostic panels, stacked from top to
bottom. Click any segment in the sidebar (or any shaded band on the
top plot) and the lower three panels redraw for it. Each panel
answers a single yes/no question; together they take seconds to
read.

\subsubsection{Panel A --- Signal $+$ segments: are the rides covered?}

This is the full-recording overview. Every detected ride is a shaded
band; the orange one is the segment currently selected.

\begin{figure}[H]
\centering
\begin{minipage}[t]{0.49\linewidth}
\centering
\includegraphics[width=\linewidth]{boutique/panels/clean_step3_overview__signal.png}
\par\smallskip
\textbf{\textcolor{green!50!black}{GOOD.}} Twelve clean rides (clean
S23 recording). Each shaded band sits over an obvious two-lobe shape;
the bare-signal stretches between bands are stationary and contain no
ride. Nothing to fix.
\end{minipage}
\hfill
\begin{minipage}[t]{0.49\linewidth}
\centering
\includegraphics[width=\linewidth]{boutique/panels/fp_step3_overview__signal.png}
\par\smallskip
\textbf{\textcolor{red!70!black}{BAD --- over-segmentation.}} Many
narrow segments; segment \#0 is a single tall spike (phone tap),
not a ride. Late segments tile over noisy hand-motion. Most should
be deleted (\S\ref{sec:boutique-fixing}, Pattern~B).
\end{minipage}
\caption{\textbf{Signal $+$ segments overview.} A clean session looks
like distinct lobe-pairs separated by quiet stretches; a bad session
either has too many bands (false positives) or stretches that
\emph{should} be shaded but aren't (missed rides, see
Figure~\ref{fig:bq-overview-damped}).}
\label{fig:bq-overview-cleanvsfp}
\end{figure}

\begin{figure}[H]
\centering
\includegraphics[width=0.97\linewidth]{boutique/panels/damped_step3_overview__signal.png}
\caption{\textbf{\textcolor{red!70!black}{BAD --- damped recording.}}
Phone-in-bag scenario. The detector found ten segments but the
recording contains many more rides --- each big black spike that is
\emph{not} shaded is a ride the detector missed because the
amplitude was below threshold. Treat this as "the detector knows it
is unsure"; verify visually and add segments by hand
(\S\ref{sec:boutique-fixing}, Pattern~A).}
\label{fig:bq-overview-damped}
\end{figure}

\subsubsection{Panel B --- Fitted trapezoid: does the template
match the ride?}

The bottom panel zooms into the selected segment and overlays the
two fitted trapezoid templates (red lobe~1, purple lobe~2) on the
raw signal. \emph{This is the single most informative panel for
deciding whether a segment is trustworthy.}

\begin{figure}[H]
\centering
\includegraphics[width=0.97\linewidth]{boutique/panels/clean_step3_overview__trap.png}

\smallskip
\textbf{\textcolor{green!50!black}{GOOD.}} The red and purple
templates sit on the smoothed envelope like a French curve on a
hand-drawn arc. Plateaus at $\pm 0.5\,$m/s$^2$ are typical of a
mid-speed elevator. Lobe~1 is positive then lobe~2 is negative
$\Rightarrow$ \emph{down} ride (the cabin decelerates after a
descent).
\caption{\textbf{Good trapezoid fit} (clean S23 recording, segment
\#5). The plateau height is in the right range
($\sim$0.5--1\,m/s$^2$); both lobes have similar magnitude;
the orange band marks the segment endpoints, which sit just outside
the lobes themselves.}
\label{fig:bq-trap-good}
\end{figure}

\begin{figure}[H]
\centering
\includegraphics[width=0.97\linewidth]{boutique/panels/fp_step3_seg0_FP__trap.png}

\smallskip
\textbf{\textcolor{red!70!black}{BAD --- amplitude wildly too high
(false positive).}} The lobes' fitted plateau is $\sim$5\,m/s$^2$
and the surrounding signal has $\pm 50\,$m/s$^2$ spikes (note the
y-axis goes to 50). This is a phone slam/drop, not a ride. The
fit \emph{looks} reasonable in shape but the magnitude is off by
$5\times$. Always cross-check the y-axis range of this panel
against the typical $\pm 1\,$m/s$^2$ for real elevators.
\caption{\textbf{Bad trapezoid fit --- spike artefact.} The pair
filter accepted a high-amplitude lobe pair around a phone-handling
event. Delete this segment.}
\label{fig:bq-trap-bad-spike}
\end{figure}

\begin{figure}[H]
\centering
\includegraphics[width=0.97\linewidth]{boutique/panels/damped_step3_seg0_full__trap.png}

\smallskip
\textbf{\textcolor{red!70!black}{BAD --- amplitude tiny (damped).}}
Lobes' plateau is $\pm 0.4\,$m/s$^2$; the lobe shapes are visible
but barely above the noise floor. The predictor will still report
a number, but its CI will be wide. Acceptable for a phone-in-pocket
recording that is your only data; otherwise prefer the barometer.
\caption{\textbf{Bad trapezoid fit --- damped.} A real ride
under-amplified by clothing/bag. Read $\Delta h$ but expect a wide
CI on this segment.}
\label{fig:bq-trap-bad-damped}
\end{figure}

\begin{figure}[H]
\centering
\includegraphics[width=0.97\linewidth]{boutique/panels/haari_step3_seg20_full__trap.png}

\smallskip
\textbf{\textcolor{red!70!black}{BAD --- segment too long (merged).}}
The orange band spans 22 seconds. Lobe~1 hits a real \emph{down}
takeoff at $t \approx 538$\,s, but the matched landing lobe is at
$t\approx 558$\,s --- with a 17-second quiet stretch in between.
That is two events labelled as one. Trim the segment to just the
first event and add a separate segment for the second
(\S\ref{sec:boutique-fixing}, Pattern~C).
\caption{\textbf{Bad trapezoid fit --- merged.} Two distinct events
inside a single segment; the trapezoid pair is forced across both.}
\label{fig:bq-trap-bad-merged}
\end{figure}

\subsubsection{Panel C --- Correlation panel: did the algorithm
find every lobe?}

Two lines: best $R^2$ for positive-going trapezoids (blue) and
negative-going ones (red). \emph{Two green dots} around the
selected segment $=$ one pair-confirmed ride. Other colours hint at
issues; the legend at the top of the panel maps each colour to a
verdict.

\begin{figure}[H]
\centering
\begin{minipage}[t]{0.49\linewidth}
\centering
\includegraphics[width=\linewidth]{boutique/panels/clean_step3_overview__corr.png}
\par\smallskip
\textbf{\textcolor{green!50!black}{GOOD.}} A green dot on the blue
peak ($t \approx 48$\,s, positive lobe) and a green dot on the red
peak ($t \approx 56$\,s, negative lobe), separated by the cruise
duration. Surrounding grey dots are sub-threshold peaks from
adjacent rides --- they are \emph{not} a problem.
\end{minipage}
\hfill
\begin{minipage}[t]{0.49\linewidth}
\centering
\includegraphics[width=\linewidth]{boutique/panels/damped_step3_seg0_full__corr.png}
\par\smallskip
\textbf{\textcolor{red!70!black}{BORDERLINE.}} Two green dots are
present but they are flanked by orange ("unpaired") and grey
("$|A|<$thr") dots. The segment is detected, but the surrounding
context says other lobes nearby almost paired and didn't quite ---
classic damped-recording symptom.
\end{minipage}
\caption{\textbf{Correlation panel.} Look for two green dots on
the selected segment. Many grey dots in between rides are normal;
many orange/purple dots inside or right next to a segment is a
hint to investigate.}
\label{fig:bq-corr-cleanvsbad}
\end{figure}

\subsubsection{Panel D --- Heatmaps: is the trapezoid shape
well-localised?}

The two heatmaps show how good $R^2$ is over the
\emph{template grid}: half-width $W$ on the y-axis, plateau
fraction $f$ on the x-axis. The white/red \textbf{X} marks the
best fit. A clean ride $=$ bright compact basin around the
\textbf{X}, with the \textbf{X} comfortably interior to the grid.

\begin{figure}[H]
\centering
\begin{minipage}[t]{0.49\linewidth}
\centering
\includegraphics[width=\linewidth]{boutique/panels/clean_step3_overview__heatmaps.png}
\par\smallskip
\textbf{\textcolor{green!50!black}{GOOD.}} Both lobes peak at the
same $(W^\star, f^\star) \approx (1.3\,\text{s}, 0.4)$ ---
interior to the grid. The bright basin is compact $\Rightarrow$
the optimum is well-localised.
\end{minipage}
\hfill
\begin{minipage}[t]{0.49\linewidth}
\centering
\includegraphics[width=\linewidth]{boutique/panels/damped_step3_seg0_full__heatmaps.png}
\par\smallskip
\textbf{\textcolor{red!70!black}{BAD --- argmax at the boundary.}}
The \textbf{X} is pinned at $f \approx 0.83$, $W$ small. This means
the true ride lies \emph{outside} the template grid --- usually
because the segment is too short, too damped, or both. Treat the
fit as approximate.
\end{minipage}
\caption{\textbf{Heatmaps.} A bright compact interior basin is
"good"; an \textbf{X} pinned to any edge of the grid is the warning
sign.}
\label{fig:bq-heatmap-cleanvsbad}
\end{figure}

\subsection{Step 3 fixes: when something is off, what to click}
\label{sec:boutique-fixing}

When a segment is bad in one of the ways above, the sidebar (left
column on step~3) gives you four primary controls:
\emph{Edit selected} (popover with start/end/type),
\emph{Delete} (trash icon),
\emph{$+$~Add segment} (insert a new 5-second \texttt{up}-typed
segment, then edit its bounds),
and \emph{Reset to detector output} (throw away every manual edit
and start over). The four common patterns:

\paragraph{Pattern A --- Missed ride.}
\textit{Symptom:} an obvious two-lobe shape on the top plot has no
shaded band over it; the correlation panel shows orange or grey
dots in that neighbourhood.\\
\textit{Fix:} \emph{$+$~Add segment} from the sidebar, then
\emph{Edit selected} and set start/end a couple of seconds outside
the visible lobes. Adjust \texttt{type} (up/down) to match.

\paragraph{Pattern B --- False positive.}
\textit{Symptom:} a shaded band over a stretch where the trapezoid
overlay (\S\ref{sec:boutique-good-vs-bad}, panel~B) is wildly
mis-scaled, or the heatmap argmax is pinned to the boundary, or
the signal in the band is just a single tall spike.\\
\textit{Fix:} select the row and click the trash icon.
Common at the very start and very end of recordings.

\paragraph{Pattern C --- Two rides merged.}
\textit{Symptom:} a single shaded band is much longer than a normal
ride, with a quiet stretch in the middle and the trapezoid lobes
fitted across the whole window
(Figure~\ref{fig:bq-trap-bad-merged}).\\
\textit{Fix:} \emph{Edit selected}, clip the end down to just past
the first lobe pair. Then \emph{$+$~Add segment} for the second
event and edit its bounds.

\paragraph{Pattern D --- Endpoints off by 1--2\,s.}
\textit{Symptom:} the orange band sits across the right ride but
the start or end is a second too early/late.\\
\textit{Fix:} \emph{Edit selected} and shift the bound. The
trapezoid overlay updates instantly, so you can see the effect.
This is harmless for predictions but worth fixing if you intend
to use the segment list downstream.

\paragraph{When in doubt, leave it alone.} The detector is
calibrated to be conservative on false positives at the cost of
occasional misses. Over-editing on borderline segments is the
most common mistake. If you cannot \emph{clearly} see a missed
ride or a false positive, accept the detector's output as-is and
move on.

\subsection{Step 4 cheat-sheet: do the two algorithms agree?}
\label{sec:boutique-step4-cheat}

The prediction step runs the Trapezoid pulse-pair algorithm
(blue bars) and ZUPT (green bars) over your final segment list.
Both should give similar $\Delta h$ on a clean recording.

\begin{figure}[H]
\centering
\begin{minipage}[t]{0.49\linewidth}
\centering
\includegraphics[width=\linewidth]{boutique/crops/clean_step4_metrics.png}
\par\smallskip
\textbf{\textcolor{green!50!black}{GOOD.}} Bars at every segment
have nearly identical heights for the two algorithms. The largest
ride (\#8, $\Delta h \approx +22$\,m) lines up. Confidence in the
report is high.
\end{minipage}
\hfill
\begin{minipage}[t]{0.49\linewidth}
\centering
\includegraphics[width=\linewidth]{boutique/haari_step4_overview.png}
\par\smallskip
\textbf{\textcolor{red!70!black}{INVESTIGATE.}} Several segments
show large blue/green disagreement. Each disagreement $>1.5$\,m
is a hint that the segment is noisy or has wrong endpoints
(\S\ref{sec:boutique-fixing}). Go back to step~3, fix that
segment, predict again.
\end{minipage}
\caption{\textbf{Step~4 bar chart.} Trapezoid (blue) and ZUPT
(green) side by side. Agreement is the eyeball check; the
detailed numbers (CI, quality, accept/reject) live below the
chart and in the wide table.}
\label{fig:bq-step4-cleanvshaari}
\end{figure}

\paragraph{Per-segment details.} Click any segment to see four
numbers per algorithm:

\begin{itemize}
  \item $\Delta h$ --- signed height (m).
  \item $\pm$CI 90\,\% --- half-width of the 90\,\% confidence
    interval. Smaller is better; below 1.5\,m is the design
    target on short rides.
  \item Quality --- a unitless score; higher is better.
  \item Accepted --- yes/no. A "no" is paired with a
    \texttt{reject\_reason} caption.
\end{itemize}

If a segment is rejected (Accepted: no), the report still includes
it but flags the row. \texttt{r2\_below\_thr} means the joint $R^2$
of the trapezoid fit was too low; \texttt{score\_above\_quantile}
means the per-segment $\sigma$ score was in the calibration
distribution's top tail. In both cases the headline number is
shown, but treat it as low-confidence.

\subsection{What you ship: the Hebrew PDF report}
\label{sec:boutique-report}

Step~5 produces a multi-page A4 PDF in Hebrew. One cover page
(metadata, summary, overview plot, per-segment table) and one page
per ride (header with $\Delta h$ / CI / quality / accept-flag,
trapezoid-overlay plot, $\pm 30$\,s correlation window). The very
last page is a Hebrew explainer of the metrics so the recipient
can read it without reading this guide. Filename:
\texttt{boutique\_report\_YYYYMMDDTHHMMSSZ.pdf}.

% =====================================================================
% Advanced section --- everything below this line is optional.
% =====================================================================

\subsection{Advanced --- only if you need to dig deeper}
\label{sec:boutique-advanced}

This subsection collects diagnostics, parameter readings, and
edge-case advice that the casual user can safely skip. It is here
for the operator who wants to sanity-check a borderline segment
before keeping or deleting it.

\subsubsection{The trapezoid parameter card}
\label{sec:boutique-paramcard}

Below the trapezoid overlay, three side-by-side cards report the
fit numerically.

\begin{figure}[H]
\centering
\includegraphics[width=0.95\linewidth]{boutique/panels/clean_step3_overview__params.png}
\caption{\textbf{Parameter card.} Lobe~1 (take-off) and lobe~2
(landing) fields ($t_c$ centre time, $A$ signed peak, $W$ half-width,
$f$ flat fraction, $R^2$ local fit), plus a joint-fit row with the
shared-shape $W^\star, f^\star$ and the joint $R^2$ / heatmap
energy. Useful sanity checks: lobe~1 $A > 0$ for an
\texttt{up} ride (sign-flipped for \texttt{down}); $|A_1|$ and
$|A_2|$ within $\sim$30\,\%; joint $R^2 > 0.85$ is the regime in
which the predictor was validated.}
\label{fig:bq-param-card}
\end{figure}

\subsubsection{The wide comparison table}

Below the bar chart, the table shows every segment with both
algorithms' $\Delta h$, CI, quality, and accept/reject side by
side. The most useful diagnostic is the column-wise difference:
$\Delta h_{\text{trap}} - \Delta h_{\text{ZUPT}}$. A
disagreement $>1.5$\,m is almost always a noisy segment or wrong
endpoints --- exactly the same patterns flagged on step~3.

\subsubsection{Choosing between Trapezoid and ZUPT}

\begin{itemize}
  \item \textbf{Trapezoid pulse-pair} (the report headline). Fits
    the seven-segment S-curve velocity profile to the joint
    two-lobe accelerometer trace. Robust on clean rides; recommended
    for the PDF report.
  \item \textbf{ZUPT.} Double-integrates the accelerometer after
    pinning the start and end velocities to zero. Robust on
    \emph{very short} rides where the trapezoid model has too few
    samples to fit; less robust on long rides where integration
    drift dominates.
\end{itemize}

The PDF uses the Trapezoid result as the headline number per ride.
The ZUPT number is shown only on step~4 as a sanity check.

\subsubsection{Noisy recordings: three regimes}

\begin{description}
  \item[Damped (phone in pocket / bag).] Symptom: signed-$R^2$
    panel mostly grey; trapezoid amplitude $<0.5$\,m/s$^2$
    (Figure~\ref{fig:bq-trap-bad-damped}); heatmap argmax pinned
    to the boundary. Add missed segments by hand if you have an
    independent ground truth (barometer, video, stopwatch); the
    predictor's $\sigma$ model widens the CI automatically.
  \item[Hand-shaky.] Symptom: many orange (unpaired) and purple
    (NMS-suppressed) dots interleaved with the green ones; bar
    chart shows one or two suspiciously short ($<$1\,m) "rides"
    around the real ones. Delete the spurious segments.
  \item[HVAC / door-slam.] Symptom: the accelerometer trace under
    a pure HVAC event is flat, but a barometer-induced
    \texttt{outside} interval can occasionally land near elevator
    landings. Spot-check segment endpoints.
\end{description}

\subsubsection{Bulk edits and reset}

The \emph{Spreadsheet editor (advanced)} expander at the bottom of
the segmentation page lets you edit the segment list as a table,
useful for ten-plus edits. The \emph{Reset to detector output}
button below the panels throws away every manual edit and re-runs
the (deterministic) detector --- the safety net if a hand-edit
goes wrong.

\subsubsection{Quick checklist before shipping a report}

A six-point eye-check that takes well under a minute and is
sufficient to clear the report for sending:

\begin{enumerate}
  \item Every shaded segment on the top plot covers a visible
    two-lobe shape, with no obvious missed lobes between segments.
  \item Every selected segment's correlation panel shows a
    \emph{two-green-dots} pattern.
  \item Every selected segment's heatmap argmax is interior to
    the grid (not pinned to a boundary).
  \item Every selected segment's trapezoid overlay sits visibly
    on the smoothed signal envelope.
  \item Both algorithms agree on $\Delta h$ within $\sim$1\,m on
    each segment, and the trapezoid algorithm marks the segment
    \emph{accepted}.
  \item CI half-width below $\sim$1.5\,m on rides shorter than
    $\sim$10\,s and below $\sim$3\,m on the longest rides.
\end{enumerate}

When all six pass, the PDF is shippable as-is. When two or more
fail on a single segment, that segment is the one to investigate
first; the rest of the report is usually unaffected.

\subsubsection{FAQ}

\paragraph{The detector returned no segments.} Either the recording
contains no rides, the rides are too short to clear the duration
floor, or the signal is damped below the amplitude threshold.
\emph{$+$~Add segment} from the sidebar to insert a manual segment
and edit its bounds.

\paragraph{Far too many segments.} Hand-shaky / walking recording.
Sort the bar chart in step~4 by quality score and delete the
low-scoring candidates.

\paragraph{Algorithms disagree by $>1.5$\,m on a segment.}
Re-clip the segment endpoints (\S\ref{sec:boutique-fixing},
Pattern~D) and predict again.

\paragraph{Quality says ``Accepted: no''.} Read the
\texttt{reject\_reason} caption. Both common reasons
(\texttt{r2\_below\_thr}, \texttt{score\_above\_quantile}) are
honest signals --- the value is reported but should be treated as
low-confidence.

\paragraph{The Download button is greyed out.} You haven't
finished step~4 yet. Go back, click \emph{Generate report},
try again.

\paragraph{Can I run the pipeline programmatically?} Yes; the
underlying \texttt{Segmenter} and \texttt{Predictor} classes
(\S\ref{sec:segmentation-detection},
\S\ref{sec:height-prediction}) are imported directly. The
Streamlit app is a thin presentation layer.
